\chapter{计算几何：精确判定与适用范围}
\section{整数系数的有界半平面交}
实现见第~\pageref{compact-IntegerHalfplanes}~页。
输入是不等式 $ax+by\le c$，三个系数均为绝对值不超过 $10^9$ 的整数。
本接口要求交集有界，只输出正面积多边形；无解、点交集和线段交集都返回空向量。
它不是一般的可行性分类接口，不可将其空返回直接解读为闭半平面交集必为空。
无界输入不在约定范围内，也不偷偷添加一个人工边框。

两条非平行边界交点用齐次坐标 $(X,Y,D)$ 表示：
\[
D=a_1b_2-b_1a_2,\quad X=c_1b_2-b_1c_2,\quad Y=a_1c_2-c_1a_2.
\]
若 $D<0$，三者同时取负，使 $D>0$。
测试 $aX+bY\le cD$ 即可判断交点位于半平面内，无需除法。
系数界保证 $|X|,|Y|,D\le2\times10^{18}$，可用 64 位保存；
相等判断的交叉乘积至多为 $4\times10^{36}$，使用 128 位。
因此本接口内的构造、排序和位置判定都是精确整数运算。

按法向量的半圆与叉积排序，比较器不使用 epsilon。
同方向平行约束按正规化后的常数项保留更紧者，避免除以法向量长度。
扫描时若新约束排除队首或队尾交点，就删除相应无效边界；
每条边界最多入队、出队一次，排序后为线性时间。
最后闭合队列，清除重复顶点；有界的非正面积结果不会输出多边形。
零法向量表示常真约束或直接矛盾，先单独处理。

测试枚举所有边界对的交点，用任意精度整数逐条验证可行性，
将可行顶点集合与返回集合比较，并检查循环顺序的严格凸性。
两套版本通过 5000 组有界随机输入、平行和退化案例、近乎平行的十亿系数，
以及十万条冗余约束的规模测试；普通与 ASan/UBSan 均通过。
将有理数坐标转换为浮点、计算面积或处理浮点输入属于另一个数值问题，尚未由此验证。
WIDA 的一般半平面交条目保留部分完成状态，浮点版本与无界分类仍需补齐。

\clearpage
\section{最近点对、Minkowski 和与三维整数判定}
最近点对与 Minkowski 和见第~\pageref{compact-GeometryExtra}~页。
最近点对返回距离平方，不开方；少于两个点时返回空 optional，重复点返回零。
分治使用递归，每层按纵坐标归并，总时间 $O(n\log n)$，额外空间 $O(n)$。
坐标绝对值上限为 $10^{12}$，作差后平方使用 128 位。

Minkowski 和输入必须是严格逆时针凸包，不重复首点，不含边上的冗余共线点；
允许循环移动起点，也允许空集、单点和线段。
输入坐标以及任意两个输入点的坐标和都须在 $[-10^{12},10^{12}]$ 内。
非退化凸包按边方向归并，时间 $O(n+m)$；退化分支枚举点和后求凸包，
时间 $O((n+m)\log(n+m))$。
返回同样不重复首点的严格凸包。

两套代码通过任意精度整数距离枚举，以及对所有点和独立执行礼品包装法的对拍；
覆盖退化、循环起点和大坐标。规模测试包含二十万点最近点对，
以及 $40001$ 顶点的抛物线凸包与自身求和，结果应为原凸包各坐标的两倍。
这些属于本地验证，不表示新增在线 AC。

三维判定见第~\pageref{compact-IntegerGeometry3D}~页。
点的三个坐标为 \texttt{x,y,z}，绝对值均不超过 $10^9$。
\texttt{orient(a,b,c,d)} 返回有向六倍四面体体积，零表示四点共面，
交换任意两个点会改变符号。\texttt{collinear} 判断三点共线，
\texttt{on\_segment} 包含端点，且支持零长度线段。
传统版线段接口名为 \texttt{On\_Segment}。

点差每个分量至多 $2\times10^9$，叉积分量至多 $8\times10^{18}$，
三重积绝对值至多 $4.8\times10^{28}$，在有符号 128 位范围内。
公开的叉积、点积辅助函数不对任意 128 位向量保证不溢出；上述保证针对点差及其叉积。
任意精度整数按 $4\times4$ 行列式的 24 个排列独立计算方向，
并用二阶子式及坐标区间验证共线与在线段上。
已通过立方体八个顶点的全部有序四元组、随机十亿级坐标、换序及退化测试，
普通与 ASan/UBSan 均通过。这里不包含三维凸包构造。
